\subsection{Residual Leakage}
\label{sec:mia}

The cryptographic guarantee of \cref{sec:security} covers the contributions, not the final model, which every
party obtains in the clear after decryption. This is not a gap a stronger mechanism could close: any protocol
that hands separate parties a shared, useful model must move knowledge between them, and a model that
classifies categories a client never saw locally necessarily encodes information about the data that taught
it those categories. The meaningful question is not whether residual leakage is zero, which it cannot be, but
how large an attack surface a protocol exposes --- and our design makes that surface as small as a one-shot
federated method can. There is no alignment phase, so nothing data-derived is released before encryption, and
the only object exposed in the clear is the single decrypted model.

Read as attack surfaces, the alternatives differ by orders of magnitude. Multi-round federated learning
publishes a fresh update every round, so a curious server or peer accumulates a long sequence of views of
each client over training~\cite{shokri2017membership,nasr2019comprehensive}. One-shot methods that rely on
differential privacy ship, in a single object, essentially the whole of a client's local knowledge --- a
synthetic copy of its data, or a noised model. HE-IFD exposes only the final aggregate: the per-client
displacements are encrypted and never appear in the clear, and because there is no alignment phase there is
no prototype channel either.

We measure the residual on that one open channel with membership-inference
attacks~\cite{yeom2018privacy,salem2019ml,carlini2022membership}, including the attacks tailored to the
federated setting~\cite{yang2023fdleaks,wang2024graddiff,gu2023ldia,jagielski2023students}, querying the
released model both as an external adversary and as a fellow client that brings its own data as a prior. The
leakage stays small: on the language task the strongest shadow-model attack sits at the level of random
guessing across heterogeneity levels, so a participant learns essentially nothing about another's membership
from the shared model beyond what any usable classifier of that data would reveal.

We do not claim the released model leaks nothing. A model that gives every client coverage of classes it
never observed must carry information about the data that supplied that coverage, and an adversary with
strong priors about a client's distribution can recover some of it; this is a property of useful shared
models, not of our protocol in particular. What the protocol guarantees is that this residual is the
\emph{entire} surface --- not multiplied across the rounds of an iterative protocol, and not handed over
wholesale as released data --- and driving the exposed surface down to that inherent floor, rather than
promising to remove a leak that cannot be removed when parties set out to build a common model, is the
privacy goal the design meets.
